\section*{Chapter \ref{ch:b2:liquid-vapour-diagrams} --- Binary Liquid--Vapour Diagrams and Distillation}

\begin{solution}{exo:b2:liquid-vapour-diagrams:1}
Benzene \qty{80.1}{\degreeCelsius}, toluene \qty{110.6}{\degreeCelsius}. The liquid
of composition 0.2 starts to boil at about \qty{102}{\degreeCelsius}; its first
vapour has composition 0.38.
\end{solution}

\begin{solution}{exo:b2:liquid-vapour-diagrams:2}
$p = 0.5 \times 1.010 + 0.5 \times 0.388 = \qty{0.699}{bar}$;
$y = 0.505/0.699 = 0.722$.
\end{solution}

\begin{solution}{exo:b2:liquid-vapour-diagrams:3}
$n_V/n = (0.30 - 0.20)/(0.45 - 0.20) = 0.40$: \qty{4.0}{mol} of vapour and
\qty{6.0}{mol} of liquid. Check: $6.0 \times 0.20 + 4.0 \times 0.45 = 3.0 = 10
\times 0.30$.
\end{solution}

\begin{solution}{exo:b2:liquid-vapour-diagrams:4}
(a) $v = 2 + 2 - 1 = 3$: $T$, $p$ and $x$ are free. (b) $v = 2 + 2 - 2 = 2$; at
fixed pressure, 1. (c) The \omterm{def:b2:liquid-vapour-diagrams:azeotrope}{azeotrope} adds the relation $x = y$: $v = 1$; at fixed
pressure 0 --- the \omterm{def:b2:liquid-vapour-diagrams:azeotrope}{azeotrope} boils at a fixed temperature, like a pure liquid,
though its composition changes with the pressure.
\end{solution}

\begin{solution}{exo:b2:liquid-vapour-diagrams:5}
At the \omterm{def:b2:liquid-vapour-diagrams:bubble-dew}{dew point} $yp/p_1^* + (1 - y)p/p_2^* = 1$. At \qty{100}{\degreeCelsius}:
$0.3 \times 1.013/1.80 + 0.7 \times 1.013/0.742 = 1.124 > 1$ (too cold); at
\qty{105}{\degreeCelsius}: $0.148 + 0.824 = 0.971 < 1$. Interpolating, $T \approx
100 + 5 \times 0.124/0.153 = \qty{104}{\degreeCelsius}$. The first drop has $x =
yp/p_1^* \approx 0.3 \times 1.013/2.0 = 0.15$.
\end{solution}

\begin{solution}{exo:b2:liquid-vapour-diagrams:6}
The first drops have the composition of the vapour over the boiling liquid,
$y = 0.71$, at \qty{92.1}{\degreeCelsius}. As benzene leaves preferentially, the
liquid grows poorer in it, its boiling point rises and the distillate grows
poorer too; the last drops are almost pure toluene. Simple distillation is a
single equilibrium stage, repeated on a liquid that keeps changing: each
fraction is a mixture.
\end{solution}

\begin{solution}{exo:b2:liquid-vapour-diagrams:7}
Feed 0.05, on the water side of the \omterm{def:b2:liquid-vapour-diagrams:azeotrope}{azeotrope}: the top delivers the \omterm{def:b2:liquid-vapour-diagrams:azeotrope}{azeotrope}
(0.88 in moles, 95~\% by mass), the boiler keeps water. Feed 0.95, on the
ethanol side: the top still delivers the \omterm{def:b2:liquid-vapour-diagrams:azeotrope}{azeotrope}, the lowest boiling point,
and the boiler keeps pure ethanol.
\end{solution}

\begin{solution}{exo:b2:liquid-vapour-diagrams:8}
The mixture boils when $0.10 + p^*_{\text{water}} = 1.013$, that is
$p^*_{\text{water}} = \qty{0.91}{bar}$, near \qty{97}{\degreeCelsius}. Mass ratio
$0.10 \times 136/(0.91 \times 18.02) = 0.83$: \qty{0.83}{g} of essence per gram
of water.
\end{solution}

\begin{solution}{exo:b2:liquid-vapour-diagrams:9}
At room temperature, two layers L$_2$ and L$_1$ (composition 0.4 lies in the
gap). On heating, the layers boil together at the \omterm{def:b2:liquid-vapour-diagrams:miscibility-gap}{heteroazeotrope}
temperature, which stays constant; the first vapour has the heteroazeotropic
composition (the dot). Since 0.4 lies on the L$_2$ side of the dot, the layer
L$_1$ is used up first; the temperature then rises, the remaining L$_2$ grows
poorer in 1 along its \omterm{def:b2:liquid-vapour-diagrams:bubble-dew}{bubble curve} and the vapour follows the \omterm{def:b2:liquid-vapour-diagrams:bubble-dew}{dew curve}, until
the point reaches the \omterm{def:b2:liquid-vapour-diagrams:bubble-dew}{dew curve}: the last drop evaporates.
\end{solution}

\begin{solution}{exo:b2:liquid-vapour-diagrams:10}
$N_{\min} = \ln(99 \times 99)/\ln 2.47 = 9.19/0.904 = 10.2$, against 6.5: going
from 95~\% to 99~\% purity at both ends costs more than half as many plates
again. Each extra factor on $x/(1 - x)$ costs the same number of plates, and
purity is measured by the impurity left, which falls only geometrically.
\end{solution}

\begin{solution}{exo:b2:liquid-vapour-diagrams:11}
The \omterm{def:b2:liquid-vapour-diagrams:binary-diagram}{isobaric diagram} alone does not give it: one needs $p_1^*$ and $p_2^*$ at
\qty{80}{\degreeCelsius}, from the Antoine equations (1.010 and
\qty{0.388}{bar}). Then $p = p_2^* + x(p_1^* - p_2^*)$ and $p = 1/(y/p_1^* + (1 -
y)/p_2^*)$. The liquid is stable at high pressure (compressing a vapour
condenses it), at low temperature on the \omterm{def:b2:liquid-vapour-diagrams:binary-diagram}{isobaric diagram}: the two diagrams
are upside down with respect to each other.
\end{solution}

\begin{solution}{exo:b2:liquid-vapour-diagrams:12}
Since $y$ increases with $x$ and $y(1 - y) > 0$, $dp/dx$ has the sign of $y - x$:
the pressure rises with $x$ where the vapour is richer in 1. For an \omterm{def:b2:chemical-potential:ideal-mixture}{ideal
mixture} $y > x$ for every $0 < x < 1$ (\cref{prop:b2:liquid-vapour-diagrams:richer}):
$p$ is strictly monotonic, has no extremum, and there is no \omterm{def:b2:liquid-vapour-diagrams:azeotrope}{azeotrope}.
\end{solution}

\begin{solution}{pb:b2:liquid-vapour-diagrams:1}
\textbf{1.} $T = B/(A - \log_{10}1.013) - C$: benzene $1203.835/4.01253 + 53.226 =
\qty{353.2}{K}$, \qty{80.1}{\degreeCelsius}; toluene $1343.943/4.07266 + 53.773 =
\qty{383.8}{K}$, \qty{110.6}{\degreeCelsius}.
\textbf{2.} At \qty{363.15}{K}: $p_1^* = \qty{1.361}{bar}$, $p_2^* = \qty{0.542}{bar}$.
\textbf{3.} $x = (1.013 - 0.542)/(1.361 - 0.542) = 0.575$; $y = 0.575 \times
1.361/1.013 = 0.773$.
\textbf{4.} $x = (1.013 - 0.742)/(1.800 - 0.742) = 0.256$; $y = 0.256 \times
1.800/1.013 = 0.455$.
\textbf{5.} \omterm{def:b2:liquid-vapour-diagrams:bubble-dew}{Bubble points} (0, 110.6), (0.256, 100), (0.575, 90), (1, 80.1); \omterm{def:b2:liquid-vapour-diagrams:bubble-dew}{dew
points} (0, 110.6), (0.455, 100), (0.773, 90), (1, 80.1) --- the lens of the
figure.
\textbf{6.} $\frac12(p_1^* + p_2^*) = 1.013$: at \qty{90}{\degreeCelsius} 0.952, at
\qty{95}{\degreeCelsius} 1.102; interpolating, $90 + 5 \times 0.061/0.150 \approx
\qty{92}{\degreeCelsius}$ (\qty{92.1}{\degreeCelsius} exactly).
\textbf{7.} $v = 2$, 1 at fixed pressure: each temperature fixes $x$ and $y$,
which is why the diagram is made of two curves.
\textbf{8.} $x = (1.013 - 0.636)/(1.569 - 0.636) = 0.404$; $y = 0.404 \times
1.569/1.013 = 0.626$.
\textbf{9.} $n_V = 100 \times (0.500 - 0.404)/(0.626 - 0.404) = \qty{43}{mol}$,
$n_L = \qty{57}{mol}$.
\textbf{10.} Vapour $43 \times 0.626 = \qty{26.9}{mol}$; liquid $57 \times 0.404
= \qty{23.0}{mol}$; total \qty{49.9}{mol} $\approx 50$.
\textbf{11.} $26.9/50 = 54~\%$ of the benzene, in a vapour of only 63~\%.
\textbf{12.} At \qty{100}{\degreeCelsius} the vapour in equilibrium has $y =
0.456 < 0.5$: the feed point lies above the \omterm{def:b2:liquid-vapour-diagrams:bubble-dew}{dew curve}, the feed is entirely
vapour and nothing is separated.
\textbf{13.} At its \omterm{def:b2:liquid-vapour-diagrams:bubble-dew}{dew point}, \qty{98.8}{\degreeCelsius} (by the method of
\cref{exo:b2:liquid-vapour-diagrams:5}).
\textbf{14.} $\alpha = p_1^*/p_2^*$; at \qty{80.1}{\degreeCelsius}, $1.013/0.390 =
2.60$.
\textbf{15.} At \qty{110.6}{\degreeCelsius}, $2.377/1.013 = 2.35$.
\textbf{16.} $\sqrt{2.60 \times 2.35} = 2.47$.
\textbf{17.} $y = xp_1^*/p$ and $1 - y = (1 - x)p_2^*/p$; divide.
\textbf{18.} Nothing is drawn off: between two plates the vapour going up and
the liquid going down carry the same amount of matter and of each component,
$V y_k = L x_{k+1}$ with $V = L$, so $x_{k+1} = y_k$.
\textbf{19.} With $r = x/(1 - x)$: $r(y_k) = \alpha\,r(x_k)$ and $x_{k+1} = y_k$, so
$r(x_{k+1}) = \alpha\,r(x_k)$ and $r(x_D) = \alpha^N r(x_B)$; $N = \ln[r(x_D)/
r(x_B)]/\ln\alpha$.
\textbf{20.} $N_{\min} = \ln(19 \times 19)/\ln 2.47 = 5.889/0.904 = 6.5$.
\textbf{21.} A column must deliver product, so it runs at a finite \omterm{def:b2:liquid-vapour-diagrams:fractional-distillation}{reflux ratio},
which needs more plates; and a real plate does not reach equilibrium (its
efficiency is below one).
\textbf{22.} Seven steps, the last ending at 0.034, below 0.05: 6.5 plates in a
continuous count, seven whole ones (the reboiler counting as one).
\textbf{23.} $x = (95/46.07)/(95/46.07 + 5/18.02) = 0.881$.
\textbf{24.} At the top, at best the \omterm{def:b2:liquid-vapour-diagrams:azeotrope}{azeotrope} (0.88 in moles); at the bottom,
water. Fifty plates do not cross the \omterm{def:b2:liquid-vapour-diagrams:azeotrope}{azeotrope}.
\textbf{25.} At total reflux, $\boldsymbol{N_{\min} \approx 6.5}$ \omterm{def:b2:liquid-vapour-diagrams:fractional-distillation}{theoretical
plates}.
\end{solution}
